\documentclass[../../../include/open-logic-section]{subfiles}

\begin{document}

\olfileid{sth}{choice}{woproblem}
\olsection{सुक्रमणाचा प्रश्न}

सुक्रमण स्वीकारतो की नाही, यावर बऱ्याच गोष्टी अवलंबून आहेत. पण या तत्त्वाची चर्चा मुख्यतः त्याच्याशी सममूल्य असलेल्या निवडीच्या स्वयंसिद्धकावर केंद्रित झाली आहे. म्हणून आता त्याकडे वळू (आणि ही सममूल्यता सिद्ध करू).

\citeyear{Cantor1883} मध्ये कँटरने सुक्रमणाच्या स्वयंसिद्धकाला पाठिंबा दर्शवला. त्याच्या मते तो “विचाराचा असा नियम आहे की जो मूलभूत, फलितांनी समृद्ध, आणि विशेषतः त्याच्या सार्वत्रिक युक्ततेमुळे उल्लेखनीय वाटतो” (\citeauthor{Potter2004} \citeyear[पृ.~243]{Potter2004} मध्ये उद्धृत). पण शेवटी या “स्वयंसिद्धकाला” पुराव्याची गरज आहे, असे कँटरचे मत झाले. गणितज्ञ समुदायाचेही तसेच मत झाले.

\citeyear{Zermelo1904} मध्ये झर्मेलोने या प्रश्नाचे “समाधान” केले. त्याचे समाधान स्पष्ट करण्यासाठी काही व्याख्या लागतील.

\begin{defn}
प्रत्येक $x \in \dom{f}$ साठी $f(x) \in x$ असेल तेव्हा आणि केवळ तेव्हाच फलन $f$ हे \emph{निवडफलन} असते. $f$ हे निवडफलन असून $\dom{f} = A \setminus \{\emptyset\}$ असेल तेव्हा आणि केवळ तेव्हाच $f$ हे \emph{$A$ साठीचे निवडफलन} आहे, असे म्हणतो.
\end{defn}

सहज अर्थाने, प्रत्येक अरिक्त संच $x \in A$ साठी, $A$ चे निवडफलन त्या $x$ मधून विशिष्ट घटक $f(x)$ \emph{निवडते}. मग निवडीचे स्वयंसिद्धक असे आहे:

\begin{axiom}[निवड]
	प्रत्येक संचाला निवडफलन असते.
\end{axiom}

निवडीवरून सुक्रमण मिळते आणि उलटही मिळते, हे झर्मेलोने दाखवले:

\begin{thm}[$\ZF$ मध्ये]\ollabel{thmwochoice}
सुक्रमण आणि निवड सममूल्य आहेत.
\end{thm}

\begin{proof}
\emph{डावीकडून उजवीकडे.} $A$ हा संचांचा संच असू द्या. युतीकरणाच्या स्वयंसिद्धकामुळे $\bigcup A$ अस्तित्वात असतो; म्हणून सुक्रमणामुळे $\bigcup A$ ला सुक्रमित करणारा $<$ आहे. आता अरिक्त सदस्यांसाठी $f(x) = \text{$x$ चा $<$-लघुतम घटक}$ असे घ्या. हे $A$ साठीचे निवडफलन आहे.

\emph{उजवीकडून डावीकडे.} $A$ निश्चित करा. रिक्त संचाचे सुक्रमण रिक्त क्रमाने मिळते; म्हणून पुढे आधारसंच अरिक्त मानू. निवडीमुळे $\Pow{A} \setminus \{\emptyset\}$ साठी निवडफलन $f$ आहे. अनंतपार पुनरावर्तन वापरून पुढील फलनाची रचना करू:
\begin{align*}
	g(0) &= f(A)\\
	g(\alpha) &=
		\begin{cases}
			\text{थांबा!{}} &\text{जर }A = \funimage{g}{\alpha}\\
			f(A \setminus \funimage{g}{\alpha}) & \text{अन्यथा}\\
		\end{cases}
\end{align*}
“थांबा!” ही सूचना अधिक लांबलचक व्याख्येचे संक्षिप्त रूप आहे. म्हणजे प्रथम $A =
\funimage{g}{\alpha}$ मिळाल्यावर प्रत्येक $\delta \geq \alpha$ साठी $g(\delta) = A$ घ्या. ही स्थिर पुढील मूल्ये आधीच्या निवडी बदलत नाहीत; पुराव्यातील एकास-एकपणाचा दावा थांबण्यापूर्वीच्या भागापुरता आहे. सुरुवातीला, याने केवळ एक \emph{पद} परिभाषित होते एवढीच खात्री आहे (\olref[sth][spine][recursion]{transrecursionschema} नंतरच्या नोंदी पाहा); पण शेवटी थांबण्यापूर्वीचा भाग एका संच-प्रांतावर मर्यादित होतो आणि त्यामुळे आपल्याला फलन मिळते, हे दाखवू.

$f$ हे निवडफलन असल्यामुळे प्रत्येक $\alpha$ साठी (थांबण्यापूर्वी) $g(\alpha) =
f(A \setminus \funimage{g}{\alpha}) \in A \setminus
\funimage{g}{\alpha}$; म्हणजे $g(\alpha) \notin \funimage{g}{\alpha}$.
म्हणून $g(\alpha) = g(\beta)$ असल्यास $g(\beta) \notin
\funimage{g}{\alpha}$, म्हणजे $\beta \notin \alpha$; आणि त्याचप्रमाणे $\alpha \notin \beta$. म्हणून त्रिपर्यायाने $\alpha = \beta$. त्यामुळे या भागावर $g$ हे !!{injective} आहे.

आता पाहा की आपण थांबतो!{}; म्हणजे $A = g[\alpha]$ असलेली एक लघुतम क्रमसूचक संख्या $\alpha$ आहे. तसे नसल्यास $g$ हे !!{injective} असल्याने प्रत्येक क्रमसूचक संख्या $\alpha$ साठी $\cardle{\alpha}{A}$ मिळेल, जे \olref[hartogs]{HartogsLemma} शी विसंगत आहे. त्यामुळे थांबण्यापूर्वीच्या भागाला मर्यादित केल्यावर $\ran{g} = A$ सुद्धा मिळते.

ही तथ्ये एकत्र केल्यावर, एखाद्या क्रमसूचक संख्येपासून $A$ पर्यंत $g$ हे !!a{bijection} आहे; येथे फक्त थांबण्यापूर्वीचा भाग घेतला आहे. आता $g$ वापरून $A$ चे सुक्रमण करता येते.
\end{proof}

म्हणून सुक्रमण आणि निवड हे एकत्रच स्वीकारले किंवा नाकारले जातात. पण प्रश्न उरतो: त्यांचा स्वीकार करावा की नकार?

\end{document}
